\documentclass{article}
\usepackage{graphicx} % Required for inserting images
\usepackage{amsmath}
\usepackage{amsfonts}
\usepackage{amssymb}


%%%%CONFIGS DE LAS MARGENES Y TAMAÑO Y POSICIÑÓN DEL TEXTO
\setlength{\textwidth}{16cm}
\setlength{\textheight}{22cm}
\voffset-20mm
\hoffset-10mm
\setlength{\parindent}{4mm}
\setlength{\parskip}{0mm plus 0.5mm minus 0.5mm}

\title{The trinity force}
\author{H. Steinhaus \and 
 H. Hahn \and
 J. Schauder \and 
 J. Sanchez \and S. Banach}

\begin{document}

\maketitle

\begin{abstract}
In this article, we present the proofs of three theorems related to the foundations of Functional Analysis. The first one deals with 
the question of whether, given a linear functional defined on a subspace can be extended to the entire space without increasing its norm.
the second one states if the operators are bounded point by point, then their norms are uniformly bounded and finally if we have a continuous linear operator
and surjective  between  Banach spaces preserve 

\end{abstract}

\section{Introdution}
Despite that Stefan Banach no study mathematic formally yours apports was a big change in the era
even more after he met Hugo Steinhaus in Kraków-Poland in 1916 where Banach begins to publish articles 




\end{document}
